\par\medskip\noindent
\textbf{\SGAInlineRefTarget{sga3:xix:paragraph:4-3}4.3. Notation.}---
If \(V\) is a vector bundle over \(S\), let \(V^\times\) denote the
open subset obtained from \(V\) by removing the zero section.  Write
the group law of \(V\) multiplicatively.  The action of
\(\mathcal O_S\) on \(V\) which defines its module structure will
then be written exponentially:
\[
  \mathcal O_S\times_S V\longrightarrow V,
  \qquad
  (x,v)\longmapsto v^x.
\]
Thus
\[
  (vv')^x=v^x(v')^x,
  \qquad
  v^{x+x'}=v^xv^{x'},
  \qquad
  v^{xx'}=(v^x)^{x'}.
\]
In particular, upon restricting the operator action to
\(\mathbb G_{m,S}\), the open subset \(V^\times\) is stable and is
therefore an object with operator group \(\mathbb G_{m,S}\).  We
shall continue to write this action as
\((z,v)\longmapsto v^z\).
\SGARefTarget{sga3:xix:paragraph:4:3}

\par\medskip\noindent
\SGAInlineRefTarget{sga3:xix:definition:4-4-1}\textbf{\textit{\SGARefLink{sga3:xix:definition:4:4:1}{Definition 4.4.1}.}}---
\textit{Let \(L\) be an invertible module on \(S\), and let
\(\mathbb W(L)\) be the associated vector bundle.  Then
\(\mathbb W(L)^\times\) is a \textup{(locally trivial)} principal
homogeneous bundle under \(\mathbb G_{m,S}\).  We write}
\[
  \Gamma(S,L)^\times=\mathbb W(L)^\times(S).
\]
\SGARefTarget{sga3:xix:definition:4:4:1}

\par\medskip\noindent
\SGAInlineRefTarget{sga3:xix:scholium:4-4-2}\textbf{\textit{\SGARefLink{sga3:xix:scholium:4:4:2}{Scholium 4.4.2}.}}---
\textit{There are natural bijections among the isomorphisms of
\(\mathcal O_S\)-modules
\(\mathcal O_S\simeq L\), the isomorphisms of
\(\mathcal O_S\)-modules
\(\mathcal O_S\simeq\mathbb W(L)\), and the
sections \(S\to\mathbb W(L)^\times\).}

The correspondence is
\[
  f\longmapsto\mathbb W(f)\longmapsto\mathbb W(f)(1).
\]
It is compatible with base change.  Thus
\(\mathbb W(L)^\times\) may be regarded as the ``scheme of
trivializations of \(\mathbb W(L)\)''.
\SGARefTarget{sga3:xix:scholium:4:4:2}

